\documentclass[10pt,a4paper]{amsart}
\usepackage{amsmath,amssymb,amsthm,graphicx}
% Retained preamble for the Stacks Project chapter snapshot in this directory.
% The chapter's own preamble.tex is a copy of the shared Stacks Project preamble; the block
% below is sliced MECHANICALLY, line for line, from it. Nothing was added, removed or
% redefined; the wrapper class, the xr-hyper/hyperref wiring and the externaldocument list
% are omitted because the assembler supplies the document class.
%
% lines taken from preamble.tex: 3, 5, 9, 11, 12, 156, 157, 158, ... (51 lines total)
%
% package coverage restored from the chapter's own preamble (2026-10-13): xy, xy-2cell

\usepackage[all]{xy}
\xyoption{2cell}
\UseAllTwocells

\newenvironment{reference}{\comment}{\endcomment}
\newenvironment{slogan}{\comment}{\endcomment}
\newenvironment{history}{\comment}{\endcomment}
\theoremstyle{plain}
\newtheorem{theorem}[subsection]{Theorem}
\newtheorem{proposition}[subsection]{Proposition}
\newtheorem{lemma}[subsection]{Lemma}
\theoremstyle{definition}
\newtheorem{definition}[subsection]{Definition}
\newtheorem{example}[subsection]{Example}
\newtheorem{exercise}[subsection]{Exercise}
\newtheorem{situation}[subsection]{Situation}
\theoremstyle{remark}
\newtheorem{remark}[subsection]{Remark}
\newtheorem{remarks}[subsection]{Remarks}
\numberwithin{equation}{subsection}
\def\lim{\mathop{\mathrm{lim}}\nolimits}
\def\colim{\mathop{\mathrm{colim}}\nolimits}
\def\Spec{\mathop{\mathrm{Spec}}}
\def\Hom{\mathop{\mathrm{Hom}}\nolimits}
\def\Ext{\mathop{\mathrm{Ext}}\nolimits}
\def\SheafHom{\mathop{\mathcal{H}\!\mathit{om}}\nolimits}
\def\SheafExt{\mathop{\mathcal{E}\!\mathit{xt}}\nolimits}
\def\Sch{\mathit{Sch}}
\def\Mor{\mathop{\mathrm{Mor}}\nolimits}
\def\Ob{\mathop{\mathrm{Ob}}\nolimits}
\def\Sh{\mathop{\mathit{Sh}}\nolimits}
\def\NL{\mathop{N\!L}\nolimits}
\def\CH{\mathop{\mathrm{CH}}\nolimits}
\def\proetale{{pro\text{-}\acute{e}tale}}
\def\etale{{\acute{e}tale}}
\def\QCoh{\mathit{QCoh}}
\def\Ker{\mathop{\mathrm{Ker}}}
\def\Im{\mathop{\mathrm{Im}}}
\def\Coker{\mathop{\mathrm{Coker}}}
\def\Coim{\mathop{\mathrm{Coim}}}
\DeclareMathSymbol{\boxtimes}{\mathbin}{AMSa}{"02}
\def\QCohstack{\mathcal{QC}\!\mathit{oh}}
\def\Cohstack{\mathcal{C}\!\mathit{oh}}
\def\Spacesstack{\mathcal{S}\!\mathit{paces}}
\def\Quotfunctor{\mathrm{Quot}}
\def\Hilbfunctor{\mathrm{Hilb}}
\def\Curvesstack{\mathcal{C}\!\mathit{urves}}
\def\Polarizedstack{\mathcal{P}\!\mathit{olarized}}
\def\Complexesstack{\mathcal{C}\!\mathit{omplexes}}
\def\Pic{\mathop{\mathrm{Pic}}\nolimits}
\def\Picardstack{\mathcal{P}\!\mathit{ic}}
\def\Picardfunctor{\mathrm{Pic}}
\def\Deformationcategory{\mathcal{D}\!\mathit{ef}}
\usepackage{cleveref}
\crefname{originallemma}{original lemma}{original lemmas}
\Crefname{originallemma}{Original lemma}{Original lemmas}
\makeatletter
\newlabel{more-algebra-definition-f-power-torsion}{{05E6}{source}{}{sourcepointer-1}{}}
\newlabel{more-algebra-definition-f-power-torsion@cref}{{[originallemma][0][]05E6}{[1][0][]source}{}{}{}}
\makeatother
\crefname{originallemma}{original lemma}{original lemmas}
\Crefname{originallemma}{Original lemma}{Original lemmas}
\makeatletter
\newlabel{perfect-lemma-affine-compare-bounded}{{06Z0}{source}{}{sourcepointer-2}{}}
\newlabel{perfect-lemma-affine-compare-bounded@cref}{{[originallemma][0][]06Z0}{[1][0][]source}{}{}{}}
\makeatother
\crefname{originallemma}{original lemma}{original lemmas}
\Crefname{originallemma}{Original lemma}{Original lemmas}
\makeatletter
\newlabel{spaces-perfect-lemma-derived-quasi-coherent-small-etale-site}{{071Q}{source}{}{sourcepointer-3}{}}
\newlabel{spaces-perfect-lemma-derived-quasi-coherent-small-etale-site@cref}{{[originallemma][0][]071Q}{[1][0][]source}{}{}{}}
\makeatother
\title{Cohomology of a sheaf supported on a finite-presentation closed subspace is power-torsion}
\author{The Stacks Project Authors}
\date{stacks-project:spaces-pushouts.tex:local-snapshot-20261012; GNU FDL 1.2 or later}
\begin{document}
\maketitle
\noindent\textit{Editorial note.} Statement span is lines 1956-1967 of sources/stacks-spaces-pushouts/pinned-original-source.tex; proof span is lines 1969-1991, the complete author proof (21 body lines).\par\medskip
\noindent\textit{Editorial note.} Navigation stubs: 3 cited-but-not-included labels are supplied as non-load-bearing pointers with their pinned tags where the snapshot tag index has them.\par\medskip
\noindent\textit{Editorial note.} extra\_packages is empty on purpose: the retained preamble already loads xy with options and a second declaration clashes in the assembler wrapper.\par\medskip
\noindent\textit{Editorial note.} Text is retained verbatim, including the author's own typos and placeholders.\par\medskip
\noindent\textit{Editorial note.} Queue index 25 of queue\_spaces-pushouts.json (pb=1217, stmt 1956-1967, proof 1969-1991).\par\medskip
\noindent\textit{Editorial note.} spaces-pushouts wave-1: fresh adjudication accept (2026-10-14); six-check dispatched pool (§2 row appended 02:4x).\par\medskip
\noindent\textit{Allowed standard prerequisites (editorial).}\begin{itemize}
\item Explicit prerequisite boundary. The cited inputs are prerequisites; the argument in the retained span is not replaced by any cited result.
\end{itemize}\medskip
\ifdefined\hypertarget\hypertarget{sourcepointer-1}{}\fi
\noindent\textit{Source reference.} Original lemma 05E6: sources/stacks-spaces-pushouts/pinned-original-source.tex, source lines 1991--1991. Cited by the retained statement or proof but not part of the collected spans; supplied only as a navigation pointer. The argument does not depend on its content.\par\medskip
\ifdefined\hypertarget\hypertarget{sourcepointer-2}{}\fi
\noindent\textit{Source reference.} Original lemma 06Z0: sources/stacks-spaces-pushouts/pinned-original-source.tex, source lines 1991--1991. Cited by the retained statement or proof but not part of the collected spans; supplied only as a navigation pointer. The argument does not depend on its content.\par\medskip
\ifdefined\hypertarget\hypertarget{sourcepointer-3}{}\fi
\noindent\textit{Source reference.} Original lemma 071Q: sources/stacks-spaces-pushouts/pinned-original-source.tex, source lines 1991--1991. Cited by the retained statement or proof but not part of the collected spans; supplied only as a navigation pointer. The argument does not depend on its content.\par\medskip
\begin{lemma}
\label{lemma-stalk-of-module-with-support}
Let $S$ be a scheme. Let $X$ be an algebraic space over $S$.
Let $i : Z \to X$ be a closed immersion of finite presentation.
Let $Q \in D_\QCoh(\mathcal{O}_X)$ be supported on $|Z|$.
Let $\overline{x}$ be a geometric point of $X$ and let
$I_{\overline{x}} \subset \mathcal{O}_{X, \overline{x}}$ be the stalk of
the ideal sheaf of $Z$. Then the cohomology modules
$H^n(Q_{\overline{x}})$ are $I_{\overline{x}}$-power torsion
(see More on Algebra, Definition
\ref{more-algebra-definition-f-power-torsion}).
\end{lemma}

\begin{proof}
Choose an affine scheme $U$ and an \'etale morphism $U \to X$ such
that $\overline{x}$ lifts to a geometric point $\overline{u}$
of $U$. Then we can replace $X$ by $U$, $Z$ by $U \times_X Z$,
$Q$ by the restriction $Q|_U$, and $\overline{x}$ by $\overline{u}$.
Thus we may assume that $X = \Spec(A)$ is affine. Let $I \subset A$
be the ideal defining $Z$. Since $i : Z \to X$ is of finite presentation,
the ideal $I = (f_1, \ldots, f_r)$ is finitely generated.
The object $Q$ comes from a complex of $A$-modules $M^\bullet$, see
Derived Categories of Spaces, Lemma
\ref{spaces-perfect-lemma-derived-quasi-coherent-small-etale-site}
and
Derived Categories of Schemes, Lemma
\ref{perfect-lemma-affine-compare-bounded}.
Since the cohomology sheaves of $Q$ are supported on $Z$
we see that the localization $M^\bullet_f$ is acyclic for each $f \in I$.
Take $x \in H^p(M^\bullet)$. By the above we can find $n_i$ such
that $f_i^{n_i} x = 0$ in $H^p(M^\bullet)$ for each $i$.
Then with $n = \sum n_i$ we see that $I^n$ annihilates $x$.
Thus $H^p(M^\bullet)$ is $I$-power torsion. Since the ring
map $A \to \mathcal{O}_{X, \overline{x}}$ is flat and since
$I_{\overline{x}} = I\mathcal{O}_{X, \overline{x}}$ we conclude.
\end{proof}

\end{document}
